\documentclass[conference]{IEEEtran}
\IEEEoverridecommandlockouts
% The preceding line is only needed to identify funding in the first footnote. If that is unneeded, please comment it out.
\usepackage{cite}
\usepackage{amsmath,amssymb,amsfonts}
\usepackage{algorithmic}
\usepackage{graphicx}
\usepackage{textcomp}
\usepackage{xcolor}
\def\BibTeX{{\rm B\kern-.05em{\sc i\kern-.025em b}\kern-.08em
    T\kern-.1667em\lower.7ex\hbox{E}\kern-.125emX}}
\begin{document}

\title{AANEC-Relaxed Trapped Surface Theorem and Stellar Mass-Gap Derivation}\author{Anonymous Research Plan Author}\maketitle

\begin{abstract}
This proposal introduces a formal theoretical framework for classical general relativity, specifically targeting spherically symmetric stellar collapse with isotropic pressure and regular matter fields. The planned contribution is a relaxed trapped-surface theorem that replaces pointwise energy conditions with the Averaged Null Energy Condition (AANEC) combined with an explicit averaged null-convergence assumption. By integrating these conditions with Misner-Sharp compactness thresholds, the work aims to establish a rigorous bound excluding stable non-black-hole remnants above a specific compactness-dependent mass scale within this restricted model class. The scope is strictly confined to formal derivations and counterexample searches, treating the observed stellar mass gap not as a direct theorem prediction but as an astrophysical consistency constraint subject to alternative stellar-physics explanations.

The expected outcome is a conditional result: if the proposed formal obligations are satisfied, the theorem will demonstrate that reaching the model-dependent compactness threshold forces trapped-surface formation, thereby excluding stable exotic remnants. This outcome remains hypothetical and unverified; it does not constitute an observed result or a definitive proof until the formal derivations and counterexample analyses are complete. The framework anticipates that the validity of this bound depends critically on the mathematical definition of the averaged focusing condition and the stability predicates for compact objects. Consequently, the anticipated contribution is a qualified theoretical boundary rather than a universal physical law, pending resolution of the identified formal gaps.
\end{abstract}
\begin{IEEEkeywords}
AANEC, trapped surface, Misner-Sharp compactness, Raychaudhuri equation, stellar mass gap, compact objects, global hyperbolicity, TOV limit
\end{IEEEkeywords}\section{Introduction}
\subsection{Theoretical and Observational Context}
The formation of black holes and the structure of the stellar mass gap are influenced by a complex interplay of general relativistic collapse and detailed stellar evolution. Recent gravitational-wave observations have enabled the inference of population properties for compact binary mergers, categorizing them into binary black hole, binary neutron star, and neutron star–black hole classes. These observations provide empirical constraints on the mass distribution of remnants. Concurrently, theoretical models suggest that the location of the lower edge of the black hole mass gap is dictated by the near-final core mass, which sets the resulting black hole mass. This edge has been found to be robust against variations in metallicity, internal mixing, and stellar wind mass loss. Furthermore, the use of non-isentropic equations of state in numerical relativity simulations demonstrates that delayed collapse to a black hole can occur, particularly in high-mass binaries, due to pressure support from temperature increases via shocks. These findings highlight the necessity of integrating precise geometric criteria for trapped-surface formation with robust astrophysical modeling.~\cite{cite_fd0f77b2a173c0ab,cite_7a407c865faafd11,cite_7a44e6730c44cb39}

\subsection{Identified Research Gap}
Despite the progress in both observational inference and numerical simulation, a critical gap remains in the formal mathematical understanding of trapped-surface formation under relaxed energy conditions. Current theoretical frameworks often rely on pointwise energy conditions, such as the Null Energy Condition (NEC) or Strong Energy Condition (SEC), which may be violated in realistic matter fields even when the Averaged Null Energy Condition (AANEC) is satisfied. There is a lack of rigorous, AANEC-relaxed trapped-surface theorems that explicitly incorporate averaged null-convergence/focusing conditions to derive a model-dependent compactness threshold. This gap prevents a clear separation between formal geometric bounds and astrophysical confounders, such as explosion energetics, fallback, or rotation, which are known to influence the observed remnant distribution. Without a qualified theorem that delineates its own boundary conditions, the connection between compactness thresholds and the exclusion of stable non-black-hole remnants remains ambiguous.~\cite{cite_7a44e6730c44cb39,cite_7a407c865faafd11}

\subsection{Proposed Contribution}
This research proposes the development of an AANEC-relaxed trapped-surface theorem and a corresponding stellar mass-gap derivation. The central contribution will be a formal proof or counterexample establishing whether AANEC, combined with an explicit averaged null-convergence/focusing condition, forces trapped-surface formation when a Misner-Sharp compactness threshold is reached. This work will define a restricted model class for spherically symmetric, non-rotating collapse with isotropic pressure and complete null generators. By deriving this formal bound, the proposal aims to clarify the geometric conditions under which stable non-black-hole remnants above a compactness-dependent mass scale are excluded. This formal contribution will be explicitly separated from astrophysical explanations, treating the observed mass gap as a consistency constraint rather than a direct theorem-derived prediction.

\subsection{Scope and Design Assumptions}
The scope of this proposal is strictly limited to classical general relativity under specific design assumptions. These include global hyperbolicity, regularity of the metric, and the satisfaction of AANEC along complete null generators without uncontrolled quantum violations. A key design assumption is the introduction of an explicit averaged null-convergence/focusing condition, which is treated as an independent geometric hypothesis rather than a consequence of AANEC alone. The formal and empirical claims will remain separate; the theorem will not attempt to quantitatively derive the observed 2-5 solar mass gap, nor will it claim to replace stellar-evolution or supernova-engine explanations. Instead, it will provide a bounded geometric criterion that can be tested against, but not conflated with, astrophysical models of collapse.

\section{Background, Survey, and Research Gap}
\subsection{Theoretical Foundations and Energy Condition Relaxation}
Classical singularity theorems rely on pointwise energy conditions, specifically the Null Energy Condition (NEC), to guarantee the focusing of null geodesics via the Raychaudhuri equation. However, quantum field theory systematically violates these pointwise conditions, as demonstrated by effects such as the Casimir energy and squeezed vacuum states. This necessitates a shift toward averaged energy conditions, such as the Achronal Averaged Null Energy Condition (AANEC), to preserve the logical chain leading to trapped-surface formation in realistic matter models. The proposed theoretical framework replaces strict pointwise assumptions with AANEC and an explicit averaged null-convergence condition, aiming to derive a trapped-surface threshold for spherical collapse. This relaxation is critical because while AANEC is robust against known quantum counterexamples, it does not automatically imply the local focusing required for classical theorems without additional geometric constraints.~\cite{cite_0ed84243cf171bc4}

\subsection{Alternative Mechanisms and Boundary Variables}
The formation of black holes versus neutron stars is governed by a complex interplay of progenitor structure and explosion dynamics, where compactness metrics often serve as better discriminators than initial mass alone. Stellar evolution models indicate that metallicity and binary interactions significantly modulate the mass reservoir available for collapse, with high-metallicity winds stripping envelopes and capping remnant masses. Furthermore, the equation of state (EOS) of dense nuclear matter introduces substantial variability; non-isentropic EOS can delay collapse through shock heating, altering the threshold for black hole formation. These astrophysical mechanisms act as confounders to any purely geometric theorem, suggesting that observed remnant distributions are not solely determined by general-relativistic stability bounds but are heavily influenced by microphysical and dynamical processes during the supernova event.~\cite{cite_7a44e6730c44cb39}

\subsection{Observational Inference and the Mass Gap}
Observational evidence from gravitational-wave astronomy and X-ray binaries reveals a distinct paucity of compact objects in the 2-5 solar mass range, known as the lower mass gap. While dynamical mass measurements in X-ray binaries suggest a narrow distribution of black hole masses, gravitational-wave detections have identified candidates, such as the secondary in GW190814, that reside within this gap. These observations challenge simple fallback models and suggest a bimodal explosion mechanism or a sharp stability threshold. However, current data cannot definitively rule out exotic compact objects (ECOs) or quark stars, as gravitational-wave signals alone may not distinguish these alternatives from black holes if their surface properties are sufficiently compact. Consequently, the mass gap serves as a consistency constraint for theoretical models rather than a direct proof of a specific trapped-surface theorem.~\cite{cite_cbbced417410bc6c,cite_07d6df5d6978a1fc}

\subsection{Research Gaps and Missing Premises}
A significant gap remains in the formal validation of AANEC-relaxed singularity theorems, specifically regarding the derivation of trapped-surface formation from averaged focusing conditions alone. The survey indicates that the precise assumptions required to bridge AANEC to a compactness bound are currently missing from the formal evidence base. Additionally, the boundary variables that strictly separate black-hole from non-black-hole outcomes in theoretical stellar evolution models remain uncovered, particularly concerning the interplay of common envelope phases and wind mass loss. This lack of a defined boundary variable prevents a direct mapping of the geometric theorem to the observed mass distribution, leaving the applicability of the theorem to realistic stellar remnants unresolved.

\subsection{Synthesis of Uncertainties and Decision Consequences}
The derived mass distributions remain sensitive to uncertainties in explosion and unbinding energies, which are not fully constrained by current progenitor models. This unresolved mechanism implies that variations in supernova energetics can populate or avoid the mass gap independently of any underlying geometric theorem. Therefore, the proposed research must treat the observed mass gap as an astrophysical consistency check rather than a falsifier of the formal theorem. Future work requires explicit mathematical formulations for the averaged null-convergence condition and the Misner-Sharp compactness parameter to test whether AANEC plus averaged focusing suffices to exclude stable non-black-hole remnants above a specific mass scale.

\section{Research Questions and Planned Contributions}
\subsection{Operational Research Questions}
This proposal addresses three operational research questions concerning the derivation of a trapped-surface threshold under relaxed energy conditions. First, what is the precise mathematical form of the explicit averaged null-convergence condition (A5) that distinguishes it from the standard Averaged Null Energy Condition (AANEC or A2)? Second, how is the Misner-Sharp compactness parameter C formally defined for the specific metric class under consideration (D4)? Third, how is the stability of non-black-hole remnants (A6) formally defined so as to distinguish it from regular black hole models such as Hayward or Bardeen? These questions delimit the scope of the formal analysis by identifying the specific geometric and physical assumptions that remain unresolved in the current design.

\subsection{Planned Theoretical Contributions}
The planned contributions focus on establishing a qualified AANEC trapped-surface bound for non-black-hole stellar remnants. The primary contribution is a formal proof or counterexample establishing whether AANEC, combined with the stated averaged null-convergence/focusing and stellar-structure assumptions, implies a trapped-surface/compactness bound. A secondary contribution is the rigorous evaluation of candidate witnesses (CE1 and CE2) to determine if they genuinely fail A5 or if A5 is ill-defined in the current context. Additionally, the proposal will assess whether numerical verification of the Raychaudhuri equation and trapped surface formation should be conducted for a discrete subset of test metrics. These contributions aim to clarify the validity domain of the mechanism replacement, specifically distinguishing the theorem's geometric constraints from astrophysical explanations for the stellar mass gap.

\subsection{Design Assumptions and Scope Boundaries}
The research design assumes that the observed 2-5 solar mass gap is not a direct consequence of the formal theorem but rather an astrophysical consistency constraint subject to alternative stellar-physics explanations, such as explosion energetics, fallback, equation-of-state uncertainties, rotation, magnetic fields, nonspherical dynamics, or binary interactions. Consequently, the formal analysis is restricted to spherically symmetric collapse of massive stars with isotropic pressure, assuming matter fields satisfy AANEC along complete null generators without uncontrolled quantum ANEC violations. The spacetime is assumed to be globally hyperbolic and sufficiently regular for trapped-surface arguments. These assumptions delimit the questions by excluding scenarios where nonspherical/rotating/magnetic dynamics dominate or where stable exotic compact objects evade the averaged-focusing assumptions.

\section{Idea Source Checkpoints and Direction Selection Audit}
\subsection{Scientific Attraction and Core Mechanism}
The selected direction, titled "AANEC-Relaxed Trapped Surface Theorem and Stellar Mass-Gap Derivation," is retained for its potential to refine the theoretical understanding of gravitational collapse by relaxing stringent pointwise energy conditions. The core mechanism proposes replacing standard pointwise Null Energy Condition (NEC) or Strong Energy Condition (SEC) assumptions with the Averaged Null Energy Condition (AANEC), supplemented by an explicit averaged null-convergence or focusing condition. This approach aims to derive a trapped-surface threshold using Raychaudhuri-type focusing and the Misner-Sharp compactness parameter. The scientific attraction lies in the possibility of establishing a rigorous geometric bound that excludes stable non-black-hole remnants above a certain compactness-dependent mass scale, thereby offering a theoretical explanation for the lower stellar mass gap that is distinct from standard astrophysical modeling.

\subsection{Audit of Source Checkpoints and Unresolved Links}
An audit of the available idea source checkpoints reveals that while the central hypothesis is consistently defined across the candidate, portfolio, and result files, critical formal elements remain unspecified. The checkpoints indicate a design assumption that AANEC plus an explicit averaged focusing condition implies trapped surface formation, yet the mathematical form of the averaged null-convergence condition (A5) and the precise definition of the Misner-Sharp compactness parameter (D4) are missing. Furthermore, the connection between the derived compactness bound and the observed 2-5 solar mass gap is treated as an astrophysical consistency constraint rather than a direct theorem-derived prediction. The current reasoning context identifies that AANEC alone may not suffice for pointwise focusing, highlighting a defect in the initial premise that requires the addition of the independent geometric hypothesis.

\subsection{Qualified Retained Direction and Explicit Exclusions}
The direction is retained with a qualified scope: it is a formal theorem/conjecture applicable to classical General Relativity under specific constraints, including spherical symmetry, isotropic pressure, global hyperbolicity, and regular matter fields. It explicitly excludes a quantitative derivation of the observed 2-5 solar mass gap as a direct consequence of the theorem, acknowledging that the gap may be governed by supernova explosion physics, fallback, equation-of-state uncertainties, rotation, magnetic fields, nonspherical dynamics, or binary interactions. The proposal does not claim to prove that all massive stars necessarily form black holes, nor does it replace stellar-evolution or supernova-engine explanations. Instead, it seeks to define the boundary conditions under which stable exotic compact objects are excluded within this restricted model class, leaving the astrophysical interpretation of the mass gap as a separate, confounder-sensitive inquiry.

\section{Problem Definition, Assumptions, and Hypotheses}
\subsection{Scope and Object of the Formal Theory}
The proposed research establishes a formal theorem or conjecture within classical General Relativity, specifically addressing spherically symmetric, non-rotating stellar collapse with isotropic pressure. The scientific object is an AANEC-relaxed trapped-surface theorem that defines a compactness bound for non-black-hole stellar remnants. This framework aims to replace pointwise energy condition assumptions with an Averaged Null Energy Condition (AANEC) combined with an explicit averaged null-convergence condition. The scope is strictly limited to the restricted model class defined by these geometric and matter-field constraints. It does not constitute a quantitative derivation of the observed 2-5 solar mass gap, nor does it claim to prove that all massive stars necessarily form black holes. Instead, the formal result is intended to serve as a boundary condition for stable equilibrium objects, with astrophysical observations treated as a separate consistency check.

\subsection{Core Hypothesis and Mechanism}
The central hypothesis posits that if spherically symmetric stellar collapse satisfies AANEC, global hyperbolicity, regular matter fields, complete null generators, and an explicit averaged null-convergence/focusing condition, then reaching a model-dependent compactness threshold forces trapped-surface formation. Consequently, stable non-black-hole remnants above a compactness-dependent mass scale are excluded within this restricted model class. The mechanism relies on deriving a trapped-surface threshold via Raychaudhuri-type focusing and Misner-Sharp compactness, then bounding stable non-black-hole remnants through TOV-like stellar-structure assumptions. This approach distinguishes itself by treating the observed stellar mass gap as an astrophysical consistency constraint rather than a direct theorem-derived prediction, acknowledging that explosion physics, fallback, and equation-of-state uncertainties may govern the actual distribution of remnant masses.

\subsection{Assumption Ledger and Geometric Constraints}
The formal derivation rests on a specific ledger of declared assumptions. A1 asserts that spacetime and matter fields satisfy the conditions of classical General Relativity for spherically symmetric collapse with isotropic pressure. A2 requires that matter fields satisfy the Averaged Null Energy Condition (AANEC) along complete null generators. A3 stipulates that spacetime is globally hyperbolic and sufficiently regular for trapped-surface arguments. A4 assumes that null generators in the relevant region are complete. Crucially, A5 introduces an explicit averaged null-convergence/focusing condition along the relevant null generators, which is treated as an independent geometric hypothesis rather than a consequence of AANEC alone. Finally, A6 models non-black-hole remnants as stable equilibrium objects obeying a bounded compactness or TOV-like limit. These assumptions collectively define the validity domain of the proposed theorem.

\subsection{Formal Definitions and Unresolved Variables}
Several formal objects are defined to structure the argument. D1 represents the mass of the collapsing matter or remnant object. D2 denotes the areal radius or radial coordinate. D3 defines the expansion scalar of a null geodesic congruence. D5 specifies the stress-energy tensor, and D6 represents the spacetime metric. D7 identifies the tangent vector to null geodesic generators. However, the rigorous derivation requires precise mathematical formulations for key variables that are currently unresolved. Specifically, the Misner-Sharp compactness parameter (D4) requires a defined formula for effective mass within radius r in the chosen metric. Furthermore, the explicit mathematical form of the averaged null-convergence condition (A5), which distinguishes it from the standard AANEC (A2), is not yet provided. These gaps are identified as needs for human input to ensure the logical integrity of the proposed proof obligations.

\subsection{Propositions and Proof Obligations}
Two primary propositions frame the theoretical contribution. P1 states that in spherically symmetric stellar collapse satisfying the specified conditions, reaching a model-dependent compactness threshold forces trapped-surface formation. P2 states that stable non-black-hole remnants above a compactness-dependent mass scale are excluded within the restricted model class defined by AANEC and averaged focusing. The validation of these propositions depends on closing specific proof obligations. PO1 addresses the derivation of negative expansion scalar via focusing. PO2 concerns the identification of trapped surfaces based on compactness thresholds. PO3 involves the consistency check against stable equilibrium limits. These obligations remain unresolved in the current draft, requiring formal derivation or counterexample search to move from candidate formalization to verified theorem status.

\section{Study Design and Methods}
\subsection{Unit of Analysis and Theoretical Protocol}
The unit of analysis is the formal proof and counterexample search within a restricted model class of classical general relativity. The protocol proposes replacing pointwise energy conditions with an averaged null energy condition (AANEC) combined with an explicit averaged null-convergence assumption. This design treats the derivation of a trapped-surface threshold via Raychaudhuri-type focusing as a planned mathematical method rather than an established result. The operational scope is limited to spherically symmetric, non-rotating collapse with isotropic pressure, ensuring that the formal claims remain distinct from empirical astrophysical observations.

\subsection{Comparators and Robustness Checks}
To assess the robustness of the proposed compactness bound, the design includes comparisons against alternative stellar-physics mechanisms. These comparators include supernova explosion energetics, fallback dynamics, equation-of-state uncertainties, rotation, magnetic fields, and binary evolution history. The method plans to evaluate whether the derived theorem remains discriminating when these confounders are introduced. This approach ensures that the formal mathematical structure is tested against physical boundary conditions, acknowledging that the observed stellar mass gap may be governed by mechanisms independent of the trapped-surface theorem.

\subsection{Counterexample and Boundary Analysis}
The study design mandates a rigorous counterexample search to define the validity domain of the proposed theorem. The planned method seeks to identify stable ultra-compact objects that satisfy AANEC and global hyperbolicity but fail to form a trapped surface due to quantum stress-energy effects or nonspherical dynamics. This analysis serves to map the boundary conditions where the averaged focusing assumption fails. By explicitly testing these limits, the design aims to restrict the theorem's applicability to cases where the added averaged null-convergence condition holds, thereby preventing overgeneralization to regimes dominated by exotic matter or incomplete null generators.

\subsection{Measurement and Calibration Requirements}
The operationalization of key variables, specifically the Misner-Sharp compactness parameter and the averaged focusing parameter, requires precise mathematical definitions that are currently unspecified. The design assumes that these parameters will be calibrated against standard stellar-structure models, but the specific formulas and measurement protocols remain undefined. Consequently, the study design cannot proceed to quantitative verification until these definitions are established. This gap necessitates a human decision on the appropriate mathematical formulation for the compactness threshold and the focusing parameter to ensure consistency with general relativistic conventions.

\subsection{Data Governance and Reproducibility}
The reproducibility of the formal derivations and the governance of any associated computational artifacts are currently unresolved. The design requires a structured plan for managing the mathematical proofs, counterexample searches, and simulation data, but the specific data management strategy has not been defined. Furthermore, the criteria for verifying the reproducibility of the theoretical results are pending human input. Without established governance protocols, the audit trail for the proposed theorem remains incomplete, limiting the ability to independently verify the logical steps of the derivation.

\section{Expected Outcome Branches and Conditional Conclusions}
\subsection{Supportive Conditional Branch}
If the prespecified analysis is consistent with the declared relation and planned controls do not favor alternative explanations, the result would support, but not prove, the conditional mechanism. This branch is expected but not observed, contingent on the formal proof establishing that AANEC plus the explicit averaged null-convergence/focusing condition forces trapped-surface formation when compactness exceeds the model-dependent threshold. Within this restricted model class, the conditional conclusion is that stable non-black-hole remnants above a compactness-dependent mass scale are excluded. This outcome does not constitute a quantitative derivation of the observed 2-5 solar mass gap, which remains subject to separate astrophysical consistency constraints.

\subsection{Heterogeneous and Competing Explanation Branch}
The relation may be conditional or heterogeneous if the prespecified analysis indicates variation across declared conditions, units, or measurement contexts. In this expected but not observed branch, no universal conclusion is warranted because the theorem fails or must be restricted if AANEC is violated, the added averaged null-convergence/focusing condition fails, null generators are incomplete, quantum stress-energy violates ANEC, nonspherical/rotating/magnetic dynamics dominate, or stable exotic compact objects evade the averaged-focusing assumptions. Furthermore, the mass-gap inference fails or becomes non-discriminating if explosion energetics, fallback, equation-of-state physics, pair-instability processes, or binary evolution dominate the observed remnant distribution. These boundary conditions necessitate predefined checks of plausible moderators to improve coverage and measurement comparability.

\subsection{Null and Contradictory Branch}
If the prespecified comparison does not support the declared relation or instead favors a declared alternative explanation, the proposed relation is not supported in this design boundary. This expected but not observed outcome implies that the theorem fails or must be restricted under the stated boundary conditions, and the mass-gap inference becomes non-discriminating if alternative stellar-physics mechanisms dominate. Absence of support in this context is not proof of absence generally. The conditional conclusion requires auditing construct validity and comparison adequacy, followed by revising the mechanism or boundary claim before another design iteration.

\subsection{Invalid or Uninterpretable Branch}
No scientific conclusion is warranted if prespecified quality-control, missingness, protocol-deviation, or validity criteria prevent interpretation. This expected but not observed branch requires resolving the identified validity or data-quality failure before repeating the design and obtaining human confirmation of measurement, sampling, and analysis prerequisites. The conditional conclusion is that the planned design did not yield interpretable evidence, leaving the formal theorem and its astrophysical consistency constraints unresolved.

\section{Risks, Limitations, and Human Review Requirements}
\subsection{Formal Verification and Mathematical Gaps}
The proposal relies on an explicit averaged null-convergence/focusing condition (A5) and a Misner-Sharp compactness parameter (D4), neither of which currently has a formal mathematical definition in the provided context. Because A5 is an independent geometric hypothesis rather than a direct consequence of the Averaged Null Energy Condition (AANEC), its precise formulation is critical for the validity of the proposed theorem. Without these definitions, it is impossible to verify whether candidate counterexamples, such as hypothetical regular black hole metrics, actually satisfy the theorem's assumptions or evade them due to undefined boundaries. This ambiguity constitutes a primary formal risk, as the logical derivation of trapped-surface formation cannot be completed or refuted until the mathematical structure of A5 and D4 is explicitly specified and reviewed.

\subsection{Astrophysical Confounders and Scope Limits}
The formal theorem is restricted to spherically symmetric, non-rotating collapse with isotropic pressure, which excludes the complex dynamics of realistic stellar evolution. Alternative mechanisms such as supernova explosion energetics, fallback, rotation, magnetic fields, and binary interactions are identified as significant confounders that could independently determine remnant masses. Evidence from gravitational-wave observations suggests that current models struggle to fully explain certain compact-object mass distributions, indicating that the observed stellar mass gap may not be a direct consequence of the proposed geometric bound. Consequently, the formal result must remain separate from empirical claims about the mass gap, as the latter requires astrophysical modeling that falls outside the scope of the restricted theorem.~\cite{cite_07d6df5d6978a1fc,cite_f634820e22d7a238,cite_789b47fb0241d9ae}

\subsection{Human Review and Verification Protocol}
Given the unresolved status of the formal reasoning plan, qualified human review is required before any further escalation or publication of results. Reviewers must assess whether the proposed witnesses for counterexamples truly satisfy all stated assumptions, particularly regarding global hyperbolicity and null generator completeness. Additionally, human input is needed to distinguish formal mathematical counterexamples from physically implausible or empirically alternative scenarios. This review process is essential to ensure that the distinction between formal claims and empirical evidence is maintained, preventing the conflation of theoretical bounds with observational data.

\section{Definitions, Propositions, and Proof Obligations}
\subsection{Formal Definitions and Symbol Roles}
The formal framework relies on a specific set of definitions and symbols to structure the argument. The spacetime metric is denoted by \$g\$, and the stress-energy tensor by \$T\$. The tangent vector to null geodesic generators is represented by \$k\$, while the expansion scalar of a null geodesic congruence is denoted by \$QWENSCI\_BACKSLASHtheta\$. The mass of the collapsing matter or remnant object is \$M\$, and the areal radius or radial coordinate is \$r\$. These symbols serve as the foundational variables for the subsequent propositions and proof obligations.

\subsection{Averaged Null Energy Condition}
\[
\int T_{ab} k^a k^b d\lambda \geq 0
\]

\subsection{The Averaged Null Energy Condition}
The Averaged Null Energy Condition (AANEC) is a central assumption in the proposed model. It requires that the integral of the stress-energy tensor \$T\_\{ab\}\$ along the tangent vector \$k\textasciicircum{}a\$ of null geodesic generators is non-negative. This condition relaxes pointwise energy conditions, allowing for local violations while maintaining an averaged constraint. The AANEC is assumed to hold along complete null generators, providing a basis for the focusing arguments that follow.

\subsection{Propositions on Trapped Surfaces and Remnant Exclusion}
Two primary propositions are proposed. Proposition P1 states that in spherically symmetric stellar collapse satisfying AANEC, global hyperbolicity, regular matter fields, complete null generators, and an explicit averaged null-convergence/focusing condition, reaching a model-dependent compactness threshold forces trapped-surface formation. Proposition P2 asserts that stable non-black-hole remnants above a compactness-dependent mass scale are excluded within this restricted model class. These propositions are candidate formalizations and remain unverified.

\subsection{Proof Obligations and Unresolved Inputs}
The derivation of the propositions depends on several proof obligations. PO1 involves verifying the averaged null-convergence condition (A5), which is currently unresolved due to a missing explicit mathematical form. PO2 requires establishing the trapped-surface identification criterion, which remains unresolved. PO3 concerns the consistency check against the stable equilibrium limit, also unresolved. These obligations highlight the critical dependencies on human input for rigorous formalization.

\subsection{Bounded Unknowns and Human Input Requirements}
The formal definition of the Misner-Sharp compactness parameter \$C\$ is missing, which prevents precise calculation of thresholds. This definition, including the formula for effective mass within radius \$r\$ in the chosen metric, is required for rigorous derivation. Additionally, the explicit mathematical form of the averaged null-convergence condition (A5) distinguishing it from A2 is not provided. These gaps are flagged as needs\_human\_input, indicating that human review is essential to resolve these foundational elements.

\subsection{Status of Formal Claims and Derivations}
All formal claims and derivations are currently in a proposed or unverified status. The forward derivation steps, including the application of the Raychaudhuri focusing equation and the trapped-surface identification criterion, are proposed but not verified. The final conclusion that stable non-black-hole remnants above a compactness-dependent mass scale are excluded remains unverified within the restricted model class. These claims are separated from empirical observations and are subject to human review and further formal development.

\section{Forward Derivation and Counterexample Search Plan}
\subsection{Lemma-by-Lemma Forward Derivation Plan}
The proposed derivation begins with the application of the Raychaudhuri focusing equation to null geodesic congruences. Under the explicit averaged null-convergence/focusing condition, the expansion scalar is proposed to become negative along the generators. This focusing is then linked to the Misner-Sharp compactness parameter to identify a threshold for trapped-surface formation. The final step involves a consistency check against the stable equilibrium limit, aiming to exclude non-black-hole remnants above this compactness-dependent mass scale.

\[
\theta < 0
\]

\subsection{Admissible Comparators and Failure Conditions}
The validity of the derivation depends on strict adherence to the model's boundary conditions. The theorem is expected to fail or require restriction if the Averaged Null Energy Condition is violated, if the added averaged null-convergence/focusing condition does not hold, or if null generators are incomplete. Additionally, the inference fails if nonspherical dynamics, rotation, magnetic fields, or stable exotic compact objects evade the averaged-focusing assumptions.

\subsection{Counterexample Search Strategy}
\begin{itemize}
\item The search plan prioritizes the identification of stable non-black-hole remnants that satisfy all stated assumptions yet fail to form a trapped surface. Candidate witnesses include regular black hole models or high-compactness metrics that may violate the specific averaged focusing constraint. The search will distinguish valid counterexamples from boundary cases where the object is either unstable or above the defined compactness threshold.
\end{itemize}

\subsection{Limitations in Formal Verification}
A critical limitation is the absence of a formal mathematical definition for the Misner-Sharp compactness parameter and the explicit form of the averaged null-convergence condition. Without these precise definitions, it is impossible to rigorously verify whether candidate witnesses satisfy all assumptions or to calculate the exact thresholds for trapped-surface formation.

\appendices
\section{Idea Source Checkpoints and Direction Selection Audit}
\subsection{Checkpoint Inventory and Audit Status}
The audit snapshot for the selected direction mechanism\_replacement is derived from three available checkpoints. These are presented as a static inventory of source files rather than a chronological sequence of development. The temporal ordering of the checkpoints is unknown, and no iteration identifiers, timestamps, or auditable diffs are present in the supplied context. The status is recorded as CHECKPOINTS\_AVAILABLE, indicating that the source records exist for review but do not constitute empirical evidence of idea evolution. The checkpoints are sourced from idea\_candidate.json, idea\_portfolio.json, and idea\_result.json, and are treated strictly as provenance anchors for the current proposal state.

\subsection{Snapshot Content and Scope Alignment}
The content of the checkpoints aligns with the central hypothesis and claim scope defined in the selected direction. The hypothesis specifies a formal theorem for classical General Relativity under spherical nonrotating collapse with isotropic pressure, AANEC, and an explicit averaged null-convergence/focusing condition. The claim scope explicitly excludes a quantitative derivation of the observed 2-5 solar mass gap, treating the gap instead as an astrophysical consistency constraint. The mechanism described in the snapshots involves replacing pointwise energy conditions with AANEC plus the averaged focusing assumption to derive a trapped-surface threshold via Raychaudhuri-type focusing and Misner-Sharp compactness. This alignment confirms that the audit snapshots support the theoretical boundaries of the current proposal without introducing unverified results.

\subsection{Targeted Gaps and Review Boundaries}
The checkpoints identify specific target gap IDs that remain open for review. These include gaps related to missing assumptions in the formal claim, unmapped boundary variables, and unresolved mechanisms regarding progenitor explosion energetics. The presence of these gaps indicates that the formal proof obligations and the connection to astrophysical observations are not yet closed. The audit does not resolve these gaps but preserves them as part of the design boundary. Consequently, the proposal remains conditional, and the identified gaps serve as markers for where further human input or formal derivation is required before the theorem or its astrophysical application can be considered complete.

\section{Variables, Symbols, and Operational Definitions}
\subsection{Geometric and Matter-Field Symbols}
The formal model is expressed through a restricted set of geometric and matter-field symbols. The spacetime metric g and the stress-energy tensor T define the gravitational and material background, while k denotes the tangent vector to null geodesic generators and theta denotes the expansion scalar of the corresponding null congruence. These symbols are introduced as candidate formalizations for the variables V1 through V5. The averaged null energy condition (AANEC) is operationally stated as the integral of T\_ab k\textasciicircum{}a k\textasciicircum{}b along null geodesics being non-negative. This baseline notation is sufficient to state the hypotheses, but it does not yet close the gap between a qualitative setup and a derivable theorem.

\subsection{Compactness, Mass, and the Threshold Gap}
The exclusion of stable non-black-hole remnants depends on the relation between mass and compactness. Mass M is defined as the mass of the collapsing matter or remnant object, and r is defined as the areal radius or radial coordinate. The Misner-Sharp compactness parameter C is the quantity intended to mediate the trapped-surface threshold. However, the specific formal definition of C, including the formula for effective mass within radius r in the chosen metric, remains unresolved. The threshold condition C >= C\_threshold is therefore proposed rather than established, and its dependence on model parameters is a human decision that must be fixed before any derivation can be tested.

\subsection{Astrophysical Confounders and Measurement Ambiguity}
Several variables that would connect the formal theorem to observed remnant distributions are not yet operationalized. Equation-of-state models and nuclear-physics parameters that determine whether a collapsing core forms a neutron star or black hole are not listed. The measurement or initial-condition specification for rotation and for magnetic fields in collapse models is undefined. The measurement method for stellar remnant mass is unspecified, and the metric for explosion energy, kinetic versus total, is not defined. These omissions mean that the confounders identified in the design, including equation-of-state uncertainty, rotation, magnetic fields, and explosion energetics, cannot presently be quantified or controlled.

\begin{itemize}
\item The geometric symbols g, T, k, and theta are candidate formalizations and remain unverified.
\item The compactness parameter C is a needs\_human\_input item; no threshold formula is supplied.
\item Rotation, magnetic fields, equation-of-state parameters, remnant-mass measurement, and explosion-energy metric are all unresolved operational definitions.
\item The observed 2-5 solar mass gap is treated as an astrophysical consistency constraint, not a theorem-derived prediction.
\end{itemize}

\section{Evidence Coverage, Unknown Items, and Review Checklist}
\subsection{Source-Bounded Evidence and Assumptions}
The proposed theoretical framework relies on the necessity of using averaged energy conditions rather than pointwise ones. Registered evidence indicates that pointwise energy conditions are systematically violated by quantum fields, which necessitates the use of weaker statements such as quantum energy inequalities or averaged energy conditions for broader validity. This supports the inclusion of AANEC as a foundational assumption for the formal theorem. Furthermore, the treatment of the observed stellar mass gap as an astrophysical consistency constraint rather than a direct theorem prediction is informed by evidence that the location of the lower edge of the black hole mass gap is dictated by near-final core mass, which sets the resulting black hole mass. These sources define the boundaries of the evidence-backed components of the proposal.~\cite{cite_0ed84243cf171bc4,cite_7a407c865faafd11}

\subsection{Critical Unknowns and Mathematical Definitions}
Several critical mathematical definitions required for the formal proof obligations remain unspecified in the current design. Specifically, the formal definition of the Misner-Sharp compactness parameter C is missing, which prevents the precise calculation of the compactness thresholds necessary to validate the theorem's conclusion. Additionally, the explicit mathematical form of the averaged null-convergence/focusing condition (A5) is absent. Without this formal definition, it is impossible to verify whether the proposed candidate counterexamples actually satisfy the averaged focusing condition or if they violate it, thereby rendering the status of the proof obligations unresolved. These gaps represent significant barriers to advancing the formal theory beyond a conjectural state.

\subsection{Review Requirements for Boundary Conditions}
The validity of the proposed theorem depends on strict adherence to its assumptions, particularly regarding stability and regularity. A review is required to distinguish formal counterexamples from empirical alternatives in the context of the proposed witnesses. This involves assessing whether a hypothetical regular high-compactness object violates global hyperbolicity or requires exotic matter to remain regular, which would invalidate it as a counterexample to the theorem's specific premises. Furthermore, the stability predicate used in the exclusion of non-black-hole remnants must be reviewed to ensure it does not circularly define stable objects as those below the compactness threshold, which would trivialize the counterexample search. These review items are essential for the rigorous verification of the formal claims.\section*{References}
\nocite{cite_07d6df5d6978a1fc,cite_0ed84243cf171bc4,cite_1e54e3c2b35110ad,cite_2576833bf7ef143a,cite_283d00782f1f9ecb,cite_357dd780a9e3a378,cite_364386a3569b6e66,cite_37ecbded9d629353,cite_42e341ff6471ce89,cite_572fec0a1e5d74d9,cite_5ea4dc0e1b288757,cite_63476e0b002da386,cite_72d7ff391b29c495,cite_789b47fb0241d9ae,cite_7a407c865faafd11,cite_7a44e6730c44cb39,cite_7ae92f06684c60ef,cite_977568cf7b8a55bf,cite_98c5e596c0c34b68,cite_bfae3a92efd1d48f,cite_c369df652181e8f7,cite_cbbced417410bc6c,cite_cc5bc07a219047b4,cite_d3016ce6d1faa8ec,cite_e228b5fd31e37646,cite_e7b42f305bc08498,cite_ecea60a2300e4892,cite_f634820e22d7a238,cite_fd0f77b2a173c0ab}
\bibliographystyle{IEEEtran}
\bibliography{references}
\end{document}
